% Options for packages loaded elsewhere
\PassOptionsToPackage{unicode}{hyperref}
\PassOptionsToPackage{hyphens}{url}
\documentclass[
  ignorenonframetext,
]{beamer}
\newif\ifbibliography
\usepackage{pgfpages}
\setbeamertemplate{caption}[numbered]
\setbeamertemplate{caption label separator}{: }
\setbeamercolor{caption name}{fg=normal text.fg}
\beamertemplatenavigationsymbolsempty
% remove section numbering
\setbeamertemplate{part page}{
  \centering
  \begin{beamercolorbox}[sep=16pt,center]{part title}
    \usebeamerfont{part title}\insertpart\par
  \end{beamercolorbox}
}
\setbeamertemplate{section page}{
  \centering
  \begin{beamercolorbox}[sep=12pt,center]{section title}
    \usebeamerfont{section title}\insertsection\par
  \end{beamercolorbox}
}
\setbeamertemplate{subsection page}{
  \centering
  \begin{beamercolorbox}[sep=8pt,center]{subsection title}
    \usebeamerfont{subsection title}\insertsubsection\par
  \end{beamercolorbox}
}
% Prevent slide breaks in the middle of a paragraph
\widowpenalties 1 10000
\raggedbottom
\AtBeginPart{
  \frame{\partpage}
}
\AtBeginSection{
  \ifbibliography
  \else
    \frame{\sectionpage}
  \fi
}
\AtBeginSubsection{
  \frame{\subsectionpage}
}
\usepackage{iftex}
\ifPDFTeX
  \usepackage[T1]{fontenc}
  \usepackage[utf8]{inputenc}
  \usepackage{textcomp} % provide euro and other symbols
\else % if luatex or xetex
  \usepackage{unicode-math} % this also loads fontspec
  \defaultfontfeatures{Scale=MatchLowercase}
  \defaultfontfeatures[\rmfamily]{Ligatures=TeX,Scale=1}
\fi
\usepackage{lmodern}
\ifPDFTeX\else
  % xetex/luatex font selection
\fi
% Use upquote if available, for straight quotes in verbatim environments
\IfFileExists{upquote.sty}{\usepackage{upquote}}{}
\IfFileExists{microtype.sty}{% use microtype if available
  \usepackage[]{microtype}
  \UseMicrotypeSet[protrusion]{basicmath} % disable protrusion for tt fonts
}{}
\makeatletter
\@ifundefined{KOMAClassName}{% if non-KOMA class
  \IfFileExists{parskip.sty}{%
    \usepackage{parskip}
  }{% else
    \setlength{\parindent}{0pt}
    \setlength{\parskip}{6pt plus 2pt minus 1pt}}
}{% if KOMA class
  \KOMAoptions{parskip=half}}
\makeatother
\setlength{\emergencystretch}{3em} % prevent overfull lines
\providecommand{\tightlist}{%
  \setlength{\itemsep}{0pt}\setlength{\parskip}{0pt}}
% Slide layout defaults for warning-clean scientific decks.
\usepackage{etoolbox}
\IfFileExists{xurl.sty}{\usepackage{xurl}}{}
\IfFileExists{seqsplit.sty}{\usepackage{seqsplit}}{\newcommand{\seqsplit}[1]{#1}}
\protected\def\breakseq#1{\seqsplit{#1}}
\protected\def\breaktt#1{\begingroup\ttfamily\seqsplit{#1}\endgroup}
\setlength{\emergencystretch}{6em}
\tolerance=5000
\hbadness=10000
\hfuzz=1pt
\setlength{\tabcolsep}{2pt}
\AtBeginEnvironment{longtable}{\tiny\renewcommand{\arraystretch}{0.86}\setlength{\tabcolsep}{1pt}}
\AtBeginEnvironment{tabular}{\tiny\renewcommand{\arraystretch}{0.86}\setlength{\tabcolsep}{1pt}}

% Natbib and cross-reference fallbacks — slides don't load natbib
% or cleveref, but combined-PDF manuscript prose may emit these
% commands. The fallback renders citations as a bracketed key list
% and cross-references as detokenized labels so slides don't fail on
% undefined control sequences or raw underscores. \providecommand is
% a no-op if packages load later.
\providecommand{\citep}[1]{[#1]}
\providecommand{\citet}[1]{#1}
\providecommand{\citealp}[1]{#1}
\providecommand{\citeauthor}[1]{#1}
\providecommand{\citeyear}[1]{#1}
\providecommand{\cref}[1]{\texttt{\detokenize{#1}}}
\providecommand{\Cref}[1]{\texttt{\detokenize{#1}}}
\usepackage{bookmark}
\IfFileExists{xurl.sty}{\usepackage{xurl}}{} % add URL line breaks if available
\urlstyle{same}
% fallback for those not using the hyperref driver hyperxmp:
\makeatletter
\@ifundefined{xmpquote}{\newcommand{\xmpquote}[1]{#1}}{}
\makeatother
\hypersetup{
  hidelinks,
  pdfcreator={LaTeX via pandoc}}

\author{\texorpdfstring{}{}}
\date{}

\begin{document}

\begin{frame}[fragile,allowframebreaks]{References}
\protect\phantomsection\label{references}
The formal bibliography (author/year entries resolved from
\breaktt{manuscript/references.bib} via the \texttt{\breakseq{[@citekey]}}
citations used throughout this manuscript) is generated by
\texttt{pandoc-crossref}/\texttt{natbib} and appears immediately below
this note. Every entry was verified this session via a live fetch
(Crossref API, DBLP, or the publisher's own DOI resolver) --- not taken
from training-data memory --- before being added to
\texttt{references.bib}; the fetch evidence is recorded in
\breaktt{ISA.md\ \#\#\ Verification}.

Annotated pointers, for readers who want the ``why this reference''
context that a bare bibliography entry doesn't carry:

\begin{itemize}
\tightlist
\item
  \textbf{Maturana \& Varela (1980), \emph{Autopoiesis and Cognition}}
  \breakseq{[@maturana\_varela\_1980]} --- the source of this project's own
  name. The book defines autopoiesis as the self-producing organization
  of living systems: a network of processes that continuously
  regenerates the components that in turn realize the network. This
  exemplar borrows the word deliberately, not decoratively ---
  \texttt{expand()} and \texttt{materialize()} are the
  ``self-producing'' processes, and \texttt{grammar.py}'s seed is the
  invariant that the network regenerates around every run.
\item
  \textbf{Claessen \& Hughes (2000), QuickCheck}
  \breakseq{[@claessen2000quickcheck]} and \textbf{MacIver et al.~(2019),
  Hypothesis} \breakseq{[@maciver2019hypothesis]} --- the
  property-based-testing lineage this project's own test suite descends
  from. \breaktt{tests/test\_property\_invariants.py} uses Hypothesis
  directly (not merely an homage) to check invariants like ``expansion
  is deterministic for any seed'' across generated inputs rather than
  hand-picked examples.
\item
  \textbf{Lamb \& Zacchiroli (2022), Reproducible Builds}
  \breakseq{[@reproducible\_builds]} --- the peer-reviewed framing for what
  \breaktt{manuscript/05\_reproducibility.md} claims operationally: that
  a build is reproducible when independent re-execution from the same
  source and environment produces bit-for-bit-identical (or
  hash-identical) output, and that this is a supply-chain-integrity
  property, not merely a convenience.
\item
  \textbf{Merkle (1987), digital signatures / hash trees}
  \breakseq{[@merkle\_tree\_provenance]} --- the theoretical basis for
  \breaktt{src/integrity.py}'s content-addressed provenance: a tree of
  hashes lets a verifier recompute and confirm the integrity of a large
  structure from its leaves up, without trusting the producer's say-so.
  This project's \breaktt{integrity\_profile:\ merkle} grammar slot is a
  direct, if much smaller-scale, application of that idea.
\end{itemize}
\end{frame}

\end{document}
